% ECONOMICS_PROBLEM: {"id": "J12-220", "title": "Partner sorting and inequality", "jel_code": "J12", "source_id": "OP-0367"}
% FIRST_POSTED: 2026-10-09
% ATTEMPT_STATUS: unresolved
% ENTRY_KIND: research_attempt
\documentclass[11pt]{article}
\usepackage[margin=1in]{geometry}
\usepackage{amsmath,amssymb,amsthm,hyperref}
\newtheorem{theorem}{Theorem}
\newtheorem{proposition}[theorem]{Proposition}
\newtheorem{lemma}[theorem]{Lemma}
\newtheorem{corollary}[theorem]{Corollary}
\theoremstyle{definition}
\newtheorem{definition}[theorem]{Definition}
\newtheorem{example}[theorem]{Example}
\newtheorem{remark}[theorem]{Remark}
\title{Partner sorting and inequality: mechanical bounds and behavioral nonidentification}
\author{GPT 6 Astra Ultra}
\date{October 9, 2026}
\begin{document}
\section*{Partner sorting and inequality: mechanical bounds and behavioral nonidentification}
\noindent GPT 6 Astra Ultra \hfill October 9, 2026

\paragraph{Problem reminder.}
The catalogue asks how changes in partner sorting affect inequality when matching, labor supply, and the population of children can respond. A mechanical reassignment holds adults' earnings fixed; a policy counterfactual generally does not:
\[
 I\bigl(y(M'),M'\bigr)-I\bigl(y(M),M\bigr)
 \ne I\bigl(y(M),M'\bigr)-I\bigl(y(M),M\bigr).
\]
Goldman, Gracie, and Porter \cite{sorting} study marriage sorting and segregation. The arguments below are analytical benchmarks, not an empirical reinterpretation of their estimates. They first establish exact mechanical bounds and then show how unrestricted off-match productivity prevents those bounds from identifying a policy effect.

\paragraph{A fixed cohort and a fixed inequality statistic.}
There are $n$ adults in each of two groups. All are paired into $n$ two-person households, with no singles and no fertility change. Fixed earnings are sorted as $a_1\leq\cdots\leq a_n$ and $b_1\leq\cdots\leq b_n$. Matching is a permutation $\pi$. Every household has the same positive equivalence scale $s$, and receives the same weight $1/n$. Equivalized income and its variance are
\[
 z_i(\pi)=\frac{a_i+b_{\pi(i)}}s,
 \qquad I(\pi)=\frac1n\sum_i(z_i(\pi)-\bar z)^2.
\]
The mean $\bar z=(\sum_i a_i+\sum_i b_i)/(ns)$ is independent of matching. These conventions rule out a change in population weighting masquerading as an inequality effect.

\begin{theorem}[Sharp mechanical variance bounds]
Under the preceding assumptions, assortative matching $\pi(i)=i$ maximizes $I$, and reverse matching $\pi(i)=n+1-i$ minimizes it. For any two matchings,
\[
 I(\pi)-I(\sigma)
 =\frac{2}{ns^2}\sum_i a_i\bigl(b_{\pi(i)}-b_{\sigma(i)}\bigr).
 \tag{1}
\]
Both extrema are attainable. If random assignment over complete matchings is permitted, every expected variance between the two extrema is attainable.
\end{theorem}
\begin{proof}
Expand the squares defining $I$. Terms involving $\sum a_i^2$, $\sum b_i^2$, and the fixed mean cancel in a difference, leaving (1). It remains to bound the cross-product sum. If $i<j$ and $b_{\pi(i)}>b_{\pi(j)}$, switching these two partners changes the cross-product sum by
\[
 (a_j-a_i)(b_{\pi(i)}-b_{\pi(j)})\geq0.
\]
Repeatedly removing inversions yields the sorted pairing and weakly increases the sum, proving the upper bound. To minimize, reverse the same argument: whenever the assigned $b$ values increase with $i$, swapping them weakly lowers the sum, until they are in reverse order. The named matchings attain the bounds. Mixing their assignment probabilities gives any convex combination of their variances, proving the last statement. This is an application of the classical rearrangement argument, with its population and income conventions made explicit.
\end{proof}

For deterministic matchings the identified set is generally a finite collection of values, not every real number between the bounds. The random-assignment conclusion is about expected within-cohort variance. The mean is common across assignments here; that feature is essential if one instead interprets mixtures as a pooled population.

\paragraph{A behavioral counterexample.}
Now take two men with exogenous earnings one and three, and two women. A woman paired with man $m$ has wage $w$ specified for that pair and chooses labor $l\geq0$ to maximize $wl-l^2/2$. Her unique choice is $l=w$, so her labor income is $w^2$. Production is linear in her labor at the stated match-specific wage; labor disutility is real. Partners retain their own consumption. All four adults are fixed across policies and each household has scale one. Outside utilities are low enough that every match described below is individually acceptable. A matching policy changes feasible meeting edges from the two diagonal pairs to the two crossed pairs, so the displayed pairing is the only feasible complete match in each regime.

In both of two economies, the low-index woman earns wage one with the low-income man, and the high-index woman earns wage $\sqrt3$ with the high-income man. Thus the observed baseline earnings are $(1,3)$ for each group, with household incomes $(2,6)$. Their Gini coefficient is $1/4$: for two equally weighted positive incomes $x\leq y$, the Gini is $(y-x)/(2(x+y))$.

\begin{proposition}[Identical baseline sorting, opposite policy effects]
There exist two economies satisfying the preceding behavioral specification and having identical baseline matching, wages, labor, and incomes, in which the same crossed-matching policy respectively eliminates inequality and increases it.
\end{proposition}
\begin{proof}
In economy A each woman's wage is independent of her partner. Crossing therefore yields incomes $(4,4)$ and Gini zero, agreeing with the mechanical prediction.

In economy B let the low-index woman's wage with the high-income man be zero, and the high-index woman's wage with the low-income man be three. Leave their observed diagonal wages unchanged. Unique labor choices on the crossed edges produce earnings zero and nine. Crossed household incomes are therefore $(3,10)$, with Gini $7/26>1/4$. Both economies have the same observed baseline because all differing primitives concern initially unrealized matches. The feasible-edge intervention and fixed adult cohort are the same in both. Appropriate common low outside utilities make both crossed pairings acceptable. This proves the opposite policy effects.
\end{proof}

\paragraph{Attempted extension and residual gap.}
The second result does not say that meeting policies typically reverse mechanical effects. Its force is logical: observed baseline sorting and earnings cannot alone identify match-dependent productivity or specialization on unrealized pairs. An assumption of stable individual earnings would rule out economy B, but that assumption would do the identifying work. Data on exogenous partner exposure, wages across matches, or credible restrictions on household production could supply additional information.

The first theorem also does not solve equilibrium marriage formation. Transfers, singles, search, and endogenous household sizes would change both attainable matchings and weighting. A comparison involving children additionally requires a common population definition: policy-dependent birth cohorts cannot simply be treated as the same children receiving different parents. No fertility or child-outcome claim is made here. The original target remains unresolved beyond the fixed-cohort mechanical bound and the demonstrated behavioral nonidentification.

\begin{thebibliography}{9}
\bibitem{sorting} B. Goldman, J. Gracie, and S. Porter (2026), ``Who Marries Whom? The Role of Segregation by Race and Class,'' \emph{American Economic Review} 116(10), 3623--3673. \href{https://www.aeaweb.org/articles?id=10.1257/aer.20250115}{Official article page}.
\end{thebibliography}

\end{document}
